\documentclass[12pt,a4paper]{article}

\usepackage[T2A]{fontenc}
\usepackage[utf8]{inputenc}
\usepackage[russian]{babel}
\usepackage{amsmath,amssymb,amsthm}
\usepackage{array}
\usepackage{enumitem}
\usepackage{listings}
\usepackage{booktabs}
\usepackage{geometry}
\geometry{margin=2cm}

\newtheorem{definition}{Определение}
\newtheorem{theorem}{Теорема}
\newtheorem{remark}{Замечание}

\lstset{
  basicstyle=\ttfamily\small,
  columns=fullflexible,
  frame=single,
  breaklines=true
}

\title{Билет №27 \\ \small Сеть Internet, доменная организация, семейство протоколов TCP/IP. (ИТМО) \\ Сеть Internet, доменная организация, семейство протоколов TCP/IP версий 4 и 6. Адресация и маршрутизация. Иерархия ISO/OSI на примере протоколов Internet. (СпБГУ)}
\author{Сериков Владислав Александрович}
\date{6 мая 2026}

\begin{document}

\maketitle
\tableofcontents
\newpage

\section{Сеть Internet: основные понятия}

\begin{definition}
\textbf{Internet} --- глобальная система взаимосвязанных компьютерных сетей, использующая семейство протоколов TCP/IP для обмена данными между узлами.
\end{definition}

Internet является не одной физической сетью, а \emph{сетью сетей}. Отдельные сети принадлежат организациям, операторам связи, университетам, компаниям и частным лицам. Их взаимодействие обеспечивается общими протоколами адресации, маршрутизации и передачи данных.

\subsection{Ключевые свойства Internet}

\begin{enumerate}
    \item \textbf{Пакетная коммутация.} Данные разбиваются на пакеты, каждый из которых передаётся независимо.
    \item \textbf{Децентрализованность.} Нет единого центра управления всей сетью.
    \item \textbf{Глобальная адресация.} Каждый узел может иметь IP-адрес, позволяющий идентифицировать его в сети.
    \item \textbf{Масштабируемость.} Иерархическая адресация и маршрутизация позволяют обслуживать миллиарды устройств.
    \item \textbf{Принцип end-to-end.} Сеть в основном отвечает за доставку пакетов, а сложная логика часто реализуется на конечных узлах.
\end{enumerate}

\begin{definition}
\textbf{Узел} --- устройство, подключённое к сети и участвующее в обмене данными: компьютер, сервер, маршрутизатор, смартфон, IoT-устройство.
\end{definition}

\begin{definition}
\textbf{Маршрутизатор} --- сетевое устройство, пересылающее IP-пакеты между сетями на основании таблицы маршрутизации.
\end{definition}

\begin{definition}
\textbf{Автономная система} --- совокупность IP-сетей и маршрутизаторов под единым административным управлением, имеющая общую политику маршрутизации.
\end{definition}

\section{Семейство протоколов TCP/IP}

\subsection{Общее устройство стека TCP/IP}

Семейство TCP/IP --- набор сетевых протоколов, используемых в Internet. Оно включает протоколы разных уровней: от передачи кадров в локальной сети до прикладных протоколов.

Обычно выделяют четыре уровня TCP/IP:

\begin{center}
\begin{tabular}{|c|p{7cm}|p{5cm}|}
\hline
\textbf{Уровень TCP/IP} & \textbf{Назначение} & \textbf{Примеры протоколов} \\
\hline
Прикладной & Сервисы для приложений & HTTP, DNS, SMTP, IMAP, SSH, FTP, DHCP \\
\hline
Транспортный & Передача данных между процессами & TCP, UDP, QUIC, SCTP \\
\hline
Межсетевой & Адресация и доставка пакетов между сетями & IPv4, IPv6, ICMP, ICMPv6, IPsec \\
\hline
Канальный / сетевого доступа & Передача кадров в пределах физической среды & Ethernet, Wi-Fi, PPP, ARP, NDP \\
\hline
\end{tabular}
\end{center}

Главная идея стека TCP/IP состоит в послойной инкапсуляции. Данные прикладного уровня помещаются в сегмент транспортного уровня, затем в IP-пакет, затем в кадр канального уровня.

\subsection{Инкапсуляция данных}

Пусть браузер отправляет HTTP-запрос серверу. Тогда данные проходят следующие стадии:

\[
\text{HTTP-сообщение}
\rightarrow
\text{TCP-сегмент}
\rightarrow
\text{IP-пакет}
\rightarrow
\text{Ethernet-кадр}.
\]

Схематически:

\[
\boxed{
\text{Ethernet header}
\left[
    \text{IP header}
    \left[
        \text{TCP header}
        \left[
            \text{HTTP data}
        \right]
    \right]
\right]
\text{Ethernet trailer}
}
\]

На принимающей стороне происходит обратный процесс --- декапсуляция.

\section{Модель ISO/OSI и её связь с Internet}

\subsection{Семь уровней ISO/OSI}

Модель ISO/OSI является эталонной моделью взаимодействия открытых систем. Она делит сетевое взаимодействие на семь уровней:

\begin{center}
\begin{tabular}{|c|c|p{8cm}|}
\hline
\textbf{Номер} & \textbf{Уровень} & \textbf{Назначение} \\
\hline
7 & Прикладной & Сетевые приложения и прикладные протоколы \\
\hline
6 & Представления & Форматы данных, кодирование, шифрование, сжатие \\
\hline
5 & Сеансовый & Управление сеансами связи \\
\hline
4 & Транспортный & Доставка данных между процессами \\
\hline
3 & Сетевой & Маршрутизация и логическая адресация \\
\hline
2 & Канальный & Передача кадров между соседними узлами \\
\hline
1 & Физический & Передача битов по среде \\
\hline
\end{tabular}
\end{center}

\subsection{Соответствие ISO/OSI и TCP/IP}

\begin{center}
\begin{tabular}{|p{4cm}|p{4cm}|p{6cm}|}
\hline
\textbf{OSI} & \textbf{TCP/IP} & \textbf{Примеры Internet-протоколов} \\
\hline
7. Прикладной & Прикладной & HTTP, DNS, SMTP, SSH, DHCP \\
\hline
6. Представления & Прикладной & TLS, MIME, JSON, ASN.1 \\
\hline
5. Сеансовый & Прикладной & RPC, NetBIOS, части TLS/SSH \\
\hline
4. Транспортный & Транспортный & TCP, UDP, QUIC, SCTP \\
\hline
3. Сетевой & Межсетевой & IPv4, IPv6, ICMP, ICMPv6, IPsec \\
\hline
2. Канальный & Сетевого доступа & Ethernet, Wi-Fi, PPP, ARP, NDP \\
\hline
1. Физический & Сетевого доступа & Витая пара, оптика, радиоканал \\
\hline
\end{tabular}
\end{center}

\subsection{Пример работы уровней на HTTP-запросе}

Пусть пользователь открывает страницу:

\[
\texttt{https://example.com/index.html}.
\]

Происходит следующая последовательность:

\begin{enumerate}
    \item \textbf{DNS} определяет IP-адрес имени \texttt{example.com}.
    \item \textbf{TCP} устанавливает соединение с сервером на порт 443.
    \item \textbf{TLS} устанавливает защищённый канал.
    \item \textbf{HTTP} передаёт запрос:
\begin{lstlisting}
GET /index.html HTTP/1.1
Host: example.com
\end{lstlisting}
    \item \textbf{IP} доставляет пакеты между сетями.
    \item \textbf{Ethernet/Wi-Fi} передаёт кадры между соседними устройствами.
\end{enumerate}

\section{IP: межсетевой уровень}

\subsection{Назначение IP}

\begin{definition}
\textbf{Internet Protocol, IP} --- протокол межсетевого уровня, обеспечивающий адресацию и доставку пакетов от отправителя к получателю через набор промежуточных сетей.
\end{definition}

IP предоставляет сервис \emph{best effort}, то есть доставку ``по возможности''. Он не гарантирует:

\begin{itemize}
    \item доставку пакета;
    \item сохранение порядка пакетов;
    \item отсутствие дублирования;
    \item отсутствие повреждения данных выше контроля заголовков.
\end{itemize}

Надёжность при необходимости реализуется протоколами более высокого уровня, например TCP.

\section{IPv4}

\subsection{Адресация IPv4}

IPv4 использует 32-битные адреса. Обычно они записываются в десятичной точечной нотации:

\[
192.168.1.10.
\]

Каждый адрес состоит из четырёх октетов:

\[
192.168.1.10 =
11000000.10101000.00000001.00001010_2.
\]

Общее число IPv4-адресов:

\[
2^{32} = 4\,294\,967\,296.
\]

\subsection{Сетевая часть и часть хоста}

IP-адрес делится на:

\begin{itemize}
    \item \textbf{префикс сети};
    \item \textbf{идентификатор узла}.
\end{itemize}

В современной записи используется CIDR-нотация:

\[
192.168.1.10/24.
\]

Здесь \texttt{/24} означает, что первые 24 бита задают сеть, а оставшиеся 8 бит --- узел.

\[
\underbrace{192.168.1}_{\text{сеть}}.\underbrace{10}_{\text{узел}}
\]

Для сети \texttt{192.168.1.0/24}:

\begin{itemize}
    \item адрес сети: \texttt{192.168.1.0};
    \item широковещательный адрес: \texttt{192.168.1.255};
    \item диапазон адресов узлов: \texttt{192.168.1.1}--\texttt{192.168.1.254}.
\end{itemize}

\subsection{Маска подсети}

Маска подсети показывает, какие биты адреса относятся к сети.

Например:

\[
/24 = 255.255.255.0.
\]

В двоичном виде:

\[
11111111.11111111.11111111.00000000.
\]

\subsection{Пример вычисления сети}

Пусть задан адрес:

\[
192.168.10.77/26.
\]

Маска \texttt{/26}:

\[
255.255.255.192.
\]

Размер блока:

\[
2^{32-26} = 2^6 = 64.
\]

Подсети в последнем октете начинаются с:

\[
0,\ 64,\ 128,\ 192.
\]

Адрес 77 попадает в диапазон 64--127, значит:

\begin{itemize}
    \item адрес сети: \texttt{192.168.10.64};
    \item широковещательный адрес: \texttt{192.168.10.127};
    \item адреса узлов: \texttt{192.168.10.65}--\texttt{192.168.10.126}.
\end{itemize}

\subsection{Частные IPv4-адреса}

Для локальных сетей выделены частные диапазоны:

\begin{center}
\begin{tabular}{|c|c|}
\hline
\textbf{Диапазон} & \textbf{CIDR} \\
\hline
10.0.0.0 -- 10.255.255.255 & 10.0.0.0/8 \\
\hline
172.16.0.0 -- 172.31.255.255 & 172.16.0.0/12 \\
\hline
192.168.0.0 -- 192.168.255.255 & 192.168.0.0/16 \\
\hline
\end{tabular}
\end{center}

Частные адреса не маршрутизируются напрямую в глобальном Internet. Для доступа наружу обычно используется NAT.

\subsection{NAT}

\begin{definition}
\textbf{NAT} --- механизм преобразования сетевых адресов, позволяющий нескольким устройствам локальной сети использовать один или несколько публичных IP-адресов.
\end{definition}

Наиболее распространённый вариант --- PAT, или NAPT, когда дополнительно преобразуются номера портов.

Пример:

\[
\begin{array}{c}
192.168.1.10:51514 \\
\downarrow \text{NAT} \\
203.0.113.5:40001
\end{array}
\]

Ответный пакет на \texttt{203.0.113.5:40001} маршрутизатор сопоставляет с внутренним узлом \texttt{192.168.1.10:51514}.

\subsection{Заголовок IPv4}

Основные поля заголовка IPv4:

\begin{center}
\begin{tabular}{|p{4cm}|p{9cm}|}
\hline
\textbf{Поле} & \textbf{Назначение} \\
\hline
Version & Версия протокола, для IPv4 равна 4 \\
\hline
IHL & Длина заголовка \\
\hline
Total Length & Полная длина IP-пакета \\
\hline
Identification & Идентификатор для фрагментации \\
\hline
Flags & Флаги фрагментации \\
\hline
Fragment Offset & Смещение фрагмента \\
\hline
TTL & Время жизни пакета \\
\hline
Protocol & Протокол верхнего уровня: TCP, UDP, ICMP \\
\hline
Header Checksum & Контрольная сумма заголовка \\
\hline
Source Address & Адрес отправителя \\
\hline
Destination Address & Адрес получателя \\
\hline
\end{tabular}
\end{center}

\subsection{TTL}

Поле TTL уменьшается на единицу каждым маршрутизатором. Если TTL становится равным нулю, пакет уничтожается. Это предотвращает бесконечную циркуляцию пакетов при ошибках маршрутизации.

\section{IPv6}

\subsection{Мотивация перехода к IPv6}

IPv6 был разработан для решения ряда проблем IPv4:

\begin{itemize}
    \item недостаток адресного пространства;
    \item необходимость упростить обработку заголовков;
    \item улучшенная поддержка автоконфигурации;
    \item лучшая поддержка расширений;
    \item отказ от обязательной фрагментации маршрутизаторами.
\end{itemize}

\subsection{Адресация IPv6}

IPv6 использует 128-битные адреса. Обычно они записываются в шестнадцатеричном виде группами по 16 бит:

\[
2001:0db8:0000:0000:0000:0000:0000:0001.
\]

Разрешается сокращать ведущие нули:

\[
2001:db8:0:0:0:0:0:1.
\]

Одна непрерывная последовательность нулевых групп может быть заменена на \texttt{::}:

\[
2001:db8::1.
\]

\subsection{Типы IPv6-адресов}

\begin{center}
\begin{tabular}{|p{4cm}|p{9cm}|}
\hline
\textbf{Тип} & \textbf{Назначение} \\
\hline
Unicast & Адрес одного интерфейса \\
\hline
Multicast & Адрес группы интерфейсов \\
\hline
Anycast & Адрес группы интерфейсов, доставка ближайшему по маршрутизации \\
\hline
\end{tabular}
\end{center}

В IPv6 нет широковещательной рассылки в стиле IPv4. Её функции выполняются multicast-адресами.

\subsection{Важные диапазоны IPv6}

\begin{center}
\begin{tabular}{|c|p{8cm}|}
\hline
\textbf{Диапазон} & \textbf{Назначение} \\
\hline
\texttt{::1/128} & Loopback-адрес \\
\hline
\texttt{fe80::/10} & Link-local-адреса \\
\hline
\texttt{fc00::/7} & Unique local addresses \\
\hline
\texttt{ff00::/8} & Multicast \\
\hline
\texttt{2000::/3} & Глобальные unicast-адреса \\
\hline
\texttt{2001:db8::/32} & Диапазон для документации \\
\hline
\end{tabular}
\end{center}

\subsection{Заголовок IPv6}

Базовый заголовок IPv6 имеет фиксированную длину 40 байт.

Основные поля:

\begin{center}
\begin{tabular}{|p{4cm}|p{9cm}|}
\hline
\textbf{Поле} & \textbf{Назначение} \\
\hline
Version & Версия протокола, для IPv6 равна 6 \\
\hline
Traffic Class & Класс трафика, используется для QoS \\
\hline
Flow Label & Метка потока \\
\hline
Payload Length & Длина полезной нагрузки \\
\hline
Next Header & Следующий заголовок: TCP, UDP, ICMPv6 или заголовок расширения \\
\hline
Hop Limit & Аналог TTL в IPv4 \\
\hline
Source Address & IPv6-адрес отправителя \\
\hline
Destination Address & IPv6-адрес получателя \\
\hline
\end{tabular}
\end{center}

\subsection{Отличия IPv4 и IPv6}

\begin{center}
\begin{tabular}{|p{4cm}|p{4cm}|p{4cm}|}
\hline
\textbf{Свойство} & \textbf{IPv4} & \textbf{IPv6} \\
\hline
Длина адреса & 32 бита & 128 бит \\
\hline
Запись адреса & Десятичная точечная & Шестнадцатеричная через двоеточия \\
\hline
Размер базового заголовка & Переменный, минимум 20 байт & Фиксированный, 40 байт \\
\hline
Контрольная сумма заголовка & Есть & Нет \\
\hline
Фрагментация & Возможна маршрутизаторами и узлами & Только конечными узлами \\
\hline
Широковещание & Есть & Нет, используется multicast \\
\hline
Автоконфигурация & Обычно DHCP & SLAAC и DHCPv6 \\
\hline
Разрешение адресов & ARP & Neighbor Discovery Protocol \\
\hline
\end{tabular}
\end{center}

\section{Адресация и маршрутизация}

\subsection{Понятие маршрута}

\begin{definition}
\textbf{Маршрут} --- запись, указывающая, куда следует переслать пакет для достижения некоторого IP-префикса.
\end{definition}

Типичная запись таблицы маршрутизации имеет вид:

\[
(\text{префикс назначения},\ \text{следующий узел},\ \text{интерфейс},\ \text{метрика}).
\]

Пример:

\begin{center}
\begin{tabular}{|c|c|c|}
\hline
\textbf{Назначение} & \textbf{Следующий узел} & \textbf{Интерфейс} \\
\hline
192.168.1.0/24 & напрямую & eth0 \\
\hline
10.0.0.0/8 & 192.168.1.1 & eth0 \\
\hline
0.0.0.0/0 & 192.168.1.254 & eth0 \\
\hline
\end{tabular}
\end{center}

Маршрут \texttt{0.0.0.0/0} называется маршрутом по умолчанию.

\subsection{Правило наибольшего совпадающего префикса}

При маршрутизации IP-пакета выбирается маршрут с наиболее длинным совпадающим префиксом.

\begin{definition}
\textbf{Наибольшее совпадение префикса} --- правило выбора маршрута, согласно которому среди всех маршрутов, чей префикс содержит адрес назначения, выбирается маршрут с максимальной длиной префикса.
\end{definition}

\begin{theorem}
Пусть таблица маршрутизации содержит набор префиксов
\[
P_1, P_2, \ldots, P_n.
\]
Для IP-адреса назначения $d$ пусть
\[
M = \{P_i \mid d \in P_i\}
\]
--- множество всех префиксов, содержащих $d$. Если $M \neq \varnothing$, то выбор префикса максимальной длины из $M$ даёт наиболее специфичный маршрут к адресу $d$.
\end{theorem}

\begin{proof}
IP-префикс длины $k$ фиксирует первые $k$ бит адреса. Чем больше $k$, тем меньше множество адресов, описываемое префиксом. Действительно, для IPv4 префикс длины $k$ содержит
\[
2^{32-k}
\]
адресов, а для IPv6:
\[
2^{128-k}.
\]
Если $k_1 > k_2$, то
\[
2^{32-k_1} < 2^{32-k_2}
\]
для IPv4 и аналогично
\[
2^{128-k_1} < 2^{128-k_2}
\]
для IPv6. Следовательно, более длинный префикс описывает меньшее множество адресов и является более специфичным.

Если адрес $d$ принадлежит нескольким префиксам, то префикс максимальной длины задаёт наименьшую среди доступных область адресов, всё ещё содержащую $d$. Поэтому такой маршрут является наиболее точным, то есть наиболее специфичным.
\end{proof}

\subsection{Алгоритм пересылки IP-пакета маршрутизатором}

\textbf{Вход:} IP-пакет с адресом назначения $d$ и таблица маршрутизации.

\textbf{Выход:} интерфейс и следующий узел либо отбрасывание пакета.

\begin{enumerate}
    \item Принять кадр на входном интерфейсе.
    \item Извлечь IP-пакет из кадра.
    \item Проверить корректность заголовка.
    \item Уменьшить TTL или Hop Limit на 1.
    \item Если TTL/Hop Limit стал равен 0, отбросить пакет и, как правило, отправить ICMP-сообщение об ошибке.
    \item Найти в таблице маршрутизации все маршруты, префиксы которых содержат адрес назначения.
    \item Выбрать маршрут с наибольшей длиной префикса.
    \item Определить следующий узел и выходной интерфейс.
    \item Инкапсулировать IP-пакет в новый канальный кадр.
    \item Передать кадр через выходной интерфейс.
\end{enumerate}

\subsection{Минимальный пример маршрутизации}

Пусть таблица маршрутизации содержит:

\[
\begin{array}{c|c}
\text{Префикс} & \text{Следующий узел} \\
\hline
0.0.0.0/0 & R_0 \\
10.0.0.0/8 & R_1 \\
10.1.0.0/16 & R_2 \\
10.1.2.0/24 & R_3
\end{array}
\]

Пакет имеет адрес назначения:

\[
10.1.2.55.
\]

Адрес принадлежит всем четырём префиксам:

\[
0.0.0.0/0,\quad 10.0.0.0/8,\quad 10.1.0.0/16,\quad 10.1.2.0/24.
\]

Самый длинный префикс --- \texttt{/24}, значит пакет будет отправлен через маршрут:

\[
10.1.2.0/24 \rightarrow R_3.
\]

\section{Внутренняя и внешняя маршрутизация}

\subsection{IGP и EGP}

Протоколы маршрутизации делятся на две большие группы:

\begin{itemize}
    \item \textbf{IGP} --- протоколы внутренней маршрутизации внутри автономной системы;
    \item \textbf{EGP} --- протоколы внешней маршрутизации между автономными системами.
\end{itemize}

Примеры IGP:

\begin{itemize}
    \item RIP;
    \item OSPF;
    \item IS-IS;
    \item EIGRP.
\end{itemize}

Главный протокол междоменной маршрутизации в Internet --- BGP.

\subsection{OSPF и алгоритм Дейкстры}

OSPF относится к протоколам состояния канала. Каждый маршрутизатор строит карту топологии сети и вычисляет кратчайшие пути до всех остальных узлов.

\subsubsection{Алгоритм Дейкстры}

\textbf{Вход:} взвешенный граф $G=(V,E)$ с неотрицательными весами рёбер, начальная вершина $s$.

\textbf{Выход:} кратчайшие расстояния $d[v]$ от $s$ до всех вершин $v$.

\begin{enumerate}
    \item Для всех вершин $v$ положить $d[v] = \infty$.
    \item Положить $d[s] = 0$.
    \item Все вершины пометить как непосещённые.
    \item Пока есть непосещённые вершины:
    \begin{enumerate}
        \item выбрать непосещённую вершину $u$ с минимальным $d[u]$;
        \item пометить $u$ как посещённую;
        \item для каждого соседа $v$ вершины $u$ выполнить релаксацию:
        \[
        \text{если } d[v] > d[u] + w(u,v),
        \text{ то } d[v] := d[u] + w(u,v).
        \]
    \end{enumerate}
\end{enumerate}

\subsubsection{Минимальный пример}

Пусть граф имеет рёбра:

\[
A-B:1,\quad A-C:4,\quad B-C:2,\quad B-D:5,\quad C-D:1.
\]

Нужно найти кратчайшие пути из $A$.

Начало:

\[
d[A]=0,\quad d[B]=\infty,\quad d[C]=\infty,\quad d[D]=\infty.
\]

После обработки $A$:

\[
d[B]=1,\quad d[C]=4.
\]

Следующая вершина $B$:

\[
d[C]=\min(4,1+2)=3,
\]
\[
d[D]=1+5=6.
\]

Следующая вершина $C$:

\[
d[D]=\min(6,3+1)=4.
\]

Итог:

\[
d[A]=0,\quad d[B]=1,\quad d[C]=3,\quad d[D]=4.
\]

Кратчайший путь из $A$ в $D$:

\[
A \rightarrow B \rightarrow C \rightarrow D.
\]

\begin{theorem}
Алгоритм Дейкстры для графа с неотрицательными весами рёбер находит кратчайшие расстояния от начальной вершины до всех достижимых вершин.
\end{theorem}

\begin{proof}
Докажем по индукции по числу посещённых вершин.

Пусть $S$ --- множество уже посещённых вершин. Инвариант алгоритма: для каждой вершины $u \in S$ значение $d[u]$ равно длине кратчайшего пути из $s$ в $u$.

База: в начале алгоритма выбирается вершина $s$, для неё $d[s]=0$. Кратчайший путь из $s$ в $s$ имеет длину 0, следовательно инвариант верен.

Переход: пусть инвариант верен для множества $S$. Алгоритм выбирает непосещённую вершину $u$ с минимальным значением $d[u]$. Нужно доказать, что $d[u]$ уже является истинным кратчайшим расстоянием.

Предположим противное: существует путь $P$ из $s$ в $u$ длины меньше $d[u]$. Так как $u \notin S$, на пути $P$ есть первая вершина $y$, не принадлежащая $S$. Пусть $x$ --- предыдущая вершина на этом пути, тогда $x \in S$.

По предположению индукции $d[x]$ равно кратчайшему расстоянию до $x$. При обработке вершины $x$ алгоритм выполнил релаксацию ребра $(x,y)$, поэтому
\[
d[y] \le d[x] + w(x,y).
\]
Так как веса рёбер неотрицательны, длина части пути от $s$ до $y$ не больше длины всего пути $P$ до $u$. Следовательно,
\[
d[y] < d[u].
\]
Но $y$ была непосещённой вершиной с расстоянием меньше $d[u]$, что противоречит выбору $u$ как непосещённой вершины с минимальным $d$. Значит, предположение неверно, и $d[u]$ является кратчайшим расстоянием.

Инвариант сохраняется. После завершения алгоритма все достижимые вершины посещены, следовательно для них найдены кратчайшие расстояния.
\end{proof}

\subsection{BGP}

\begin{definition}
\textbf{BGP} --- протокол междоменной маршрутизации, используемый для обмена маршрутной информацией между автономными системами.
\end{definition}

BGP является протоколом типа path-vector. Он передаёт не только информацию о достижимом префиксе, но и путь через автономные системы.

Пример BGP-маршрута:

\[
203.0.113.0/24 \rightarrow AS64500,\ AS64496,\ AS64497.
\]

Здесь \texttt{AS\_PATH} показывает последовательность автономных систем, через которые проходит маршрут.

Основные атрибуты BGP:

\begin{itemize}
    \item \texttt{AS\_PATH} --- путь через автономные системы;
    \item \texttt{NEXT\_HOP} --- следующий маршрутизатор;
    \item \texttt{LOCAL\_PREF} --- локальное предпочтение маршрута;
    \item \texttt{MED} --- метрика для выбора входа в соседнюю автономную систему;
    \item \texttt{COMMUNITY} --- метки для политики маршрутизации.
\end{itemize}

BGP выбирает маршрут не просто по минимальному числу переходов, а по политике. Это важно, потому что Internet состоит из независимых организаций с собственными экономическими и техническими интересами.

\section{ARP и NDP}

\subsection{ARP в IPv4}

\begin{definition}
\textbf{ARP} --- протокол определения канального адреса по IPv4-адресу в локальной сети.
\end{definition}

Если узел \texttt{192.168.1.10} хочет отправить пакет узлу \texttt{192.168.1.20} в той же сети, ему нужен MAC-адрес получателя.

Алгоритм:

\begin{enumerate}
    \item Узел проверяет ARP-кэш.
    \item Если MAC-адрес неизвестен, отправляется широковещательный ARP-запрос:
\begin{lstlisting}
Who has 192.168.1.20? Tell 192.168.1.10
\end{lstlisting}
    \item Узел с адресом \texttt{192.168.1.20} отвечает:
\begin{lstlisting}
192.168.1.20 is at aa:bb:cc:dd:ee:ff
\end{lstlisting}
    \item Отправитель сохраняет соответствие в ARP-кэше.
    \item Далее IP-пакет инкапсулируется в Ethernet-кадр с найденным MAC-адресом.
\end{enumerate}

\subsection{NDP в IPv6}

В IPv6 вместо ARP используется Neighbor Discovery Protocol. Он основан на ICMPv6 и выполняет функции:

\begin{itemize}
    \item обнаружение соседей;
    \item определение канальных адресов;
    \item обнаружение маршрутизаторов;
    \item обнаружение дублирующихся адресов;
    \item поддержка автоконфигурации.
\end{itemize}

\section{ICMP и ICMPv6}

\begin{definition}
\textbf{ICMP} --- протокол управляющих сообщений, используемый для диагностики и передачи ошибок на сетевом уровне.
\end{definition}

Примеры ICMP-сообщений:

\begin{itemize}
    \item Echo Request;
    \item Echo Reply;
    \item Destination Unreachable;
    \item Time Exceeded.
\end{itemize}

Команда \texttt{ping} использует ICMP Echo Request и Echo Reply.

Команда \texttt{traceroute} использует уменьшение TTL/Hop Limit и получение сообщений Time Exceeded от промежуточных маршрутизаторов.

\section{Транспортный уровень}

\subsection{Порты}

IP-адрес идентифицирует узел, а порт идентифицирует процесс или сетевую службу на узле.

Сокет обычно задаётся парой:

\[
(\text{IP-адрес},\ \text{порт}).
\]

TCP-соединение определяется пятёркой:

\[
(\text{IP}_{src},\ \text{port}_{src},\ \text{IP}_{dst},\ \text{port}_{dst},\ \text{protocol}).
\]

Примеры известных портов:

\begin{center}
\begin{tabular}{|c|c|c|}
\hline
\textbf{Порт} & \textbf{Протокол} & \textbf{Назначение} \\
\hline
53 & TCP/UDP & DNS \\
\hline
80 & TCP & HTTP \\
\hline
443 & TCP/UDP & HTTPS, HTTP/3 поверх QUIC \\
\hline
22 & TCP & SSH \\
\hline
25 & TCP & SMTP \\
\hline
\end{tabular}
\end{center}

\subsection{UDP}

\begin{definition}
\textbf{UDP} --- транспортный протокол без установления соединения, предоставляющий передачу отдельных датаграмм между процессами.
\end{definition}

Свойства UDP:

\begin{itemize}
    \item нет установления соединения;
    \item нет гарантированной доставки;
    \item нет контроля порядка;
    \item малая служебная нагрузка;
    \item подходит для DNS, VoIP, потокового видео, игр, QUIC.
\end{itemize}

Заголовок UDP содержит только четыре поля:

\begin{center}
\begin{tabular}{|c|p{8cm}|}
\hline
\textbf{Поле} & \textbf{Назначение} \\
\hline
Source Port & Порт отправителя \\
\hline
Destination Port & Порт получателя \\
\hline
Length & Длина UDP-датаграммы \\
\hline
Checksum & Контрольная сумма \\
\hline
\end{tabular}
\end{center}

\subsection{TCP}

\begin{definition}
\textbf{TCP} --- транспортный протокол с установлением соединения, обеспечивающий надёжный упорядоченный байтовый поток между двумя процессами.
\end{definition}

TCP обеспечивает:

\begin{itemize}
    \item установление соединения;
    \item надёжную доставку;
    \item сохранение порядка байтов;
    \item контроль потока;
    \item контроль перегрузки;
    \item обнаружение потерь и повторную передачу.
\end{itemize}

\subsection{Трёхстороннее рукопожатие TCP}

Перед передачей данных TCP устанавливает соединение.

\textbf{Шаги:}

\begin{enumerate}
    \item Клиент отправляет сегмент \texttt{SYN}:
\[
C \rightarrow S: SYN,\ seq=x.
\]
    \item Сервер отвечает \texttt{SYN+ACK}:
\[
S \rightarrow C: SYN,\ ACK,\ seq=y,\ ack=x+1.
\]
    \item Клиент подтверждает:
\[
C \rightarrow S: ACK,\ seq=x+1,\ ack=y+1.
\]
\end{enumerate}

Схема:

\[
\begin{array}{ccc}
\text{Клиент} & & \text{Сервер} \\
SYN, seq=x & \longrightarrow & \\
& \longleftarrow & SYN+ACK, seq=y, ack=x+1 \\
ACK, ack=y+1 & \longrightarrow & \\
\end{array}
\]

После этого соединение считается установленным.

\subsection{Завершение TCP-соединения}

Обычно используется обмен сегментами \texttt{FIN} и \texttt{ACK}:

\[
\begin{array}{ccc}
A & & B \\
FIN & \longrightarrow & \\
& \longleftarrow & ACK \\
& \longleftarrow & FIN \\
ACK & \longrightarrow & \\
\end{array}
\]

\subsection{Скользящее окно TCP}

TCP использует механизм скользящего окна для контроля потока.

Пусть отправитель может иметь неподтверждёнными не более $W$ байтов. Тогда он может отправлять новые данные, пока количество неподтверждённых байтов меньше $W$.

Пример:

\[
W=4000.
\]

Отправитель передал байты с номерами:

\[
1,\ldots,4000.
\]

Пока подтверждения нет, новые байты отправлять нельзя. Если получено подтверждение:

\[
ACK=2001,
\]

это означает, что байты до 2000 включительно получены. Окно сдвигается, и можно отправить ещё 2000 байтов.

\section{DNS и доменная организация Internet}

\subsection{Назначение DNS}

\begin{definition}
\textbf{DNS} --- распределённая иерархическая система доменных имён, сопоставляющая человекочитаемые имена с сетевыми данными, прежде всего IP-адресами.
\end{definition}

Например:

\[
\texttt{www.example.com} \rightarrow 93.184.216.34.
\]

DNS нужен потому, что людям удобнее запоминать имена, а сетевые протоколы используют IP-адреса.

\subsection{Иерархия доменных имён}

Доменное имя имеет древовидную структуру. Корень обозначается точкой:

\[
.
\]

Пример полного доменного имени:

\[
\texttt{www.example.com.}
\]

Его части:

\[
\underbrace{\texttt{www}}_{\text{узел или поддомен}}.
\underbrace{\texttt{example}}_{\text{домен второго уровня}}.
\underbrace{\texttt{com}}_{\text{домен верхнего уровня}}.
\underbrace{\texttt{.}}_{\text{корень}}
\]

Иерархия:

\[
\texttt{.}
\rightarrow
\texttt{com}
\rightarrow
\texttt{example.com}
\rightarrow
\texttt{www.example.com}.
\]

\subsection{Зоны DNS}

\begin{definition}
\textbf{DNS-зона} --- административно управляемая часть пространства доменных имён, за которую отвечает определённый набор авторитативных DNS-серверов.
\end{definition}

Домен и зона не всегда совпадают. Например, домен \texttt{example.com} может делегировать поддомен \texttt{sub.example.com} другим DNS-серверам. Тогда \texttt{sub.example.com} станет отдельной зоной.

\subsection{Типы DNS-серверов}

\begin{itemize}
    \item \textbf{Корневые серверы} --- знают серверы доменов верхнего уровня.
    \item \textbf{TLD-серверы} --- обслуживают домены верхнего уровня: \texttt{com}, \texttt{org}, \texttt{ru} и др.
    \item \textbf{Авторитативные серверы} --- хранят окончательные записи для конкретных зон.
    \item \textbf{Рекурсивные резолверы} --- выполняют поиск имени по поручению клиента.
\end{itemize}

\subsection{Основные типы DNS-записей}

\begin{center}
\begin{tabular}{|c|p{9cm}|}
\hline
\textbf{Тип} & \textbf{Назначение} \\
\hline
A & IPv4-адрес имени \\
\hline
AAAA & IPv6-адрес имени \\
\hline
NS & Авторитативный сервер зоны \\
\hline
CNAME & Каноническое имя, псевдоним \\
\hline
MX & Почтовые серверы домена \\
\hline
TXT & Текстовые данные, часто SPF/DKIM/служебные записи \\
\hline
SOA & Начальная запись зоны, параметры зоны \\
\hline
PTR & Обратное отображение IP-адреса в имя \\
\hline
SRV & Сервисная запись с указанием порта и хоста \\
\hline
CAA & Разрешённые центры сертификации для домена \\
\hline
\end{tabular}
\end{center}

\subsection{Рекурсивное разрешение доменного имени}

Пусть клиенту нужен IP-адрес имени:

\[
\texttt{www.example.com}.
\]

\textbf{Алгоритм:}

\begin{enumerate}
    \item Клиент обращается к рекурсивному DNS-резолверу.
    \item Резолвер проверяет свой кэш.
    \item Если ответа нет, резолвер обращается к корневому серверу.
    \item Корневой сервер сообщает адреса TLD-серверов для \texttt{com}.
    \item Резолвер обращается к TLD-серверу \texttt{com}.
    \item TLD-сервер сообщает авторитативные серверы зоны \texttt{example.com}.
    \item Резолвер обращается к авторитативному серверу \texttt{example.com}.
    \item Авторитативный сервер возвращает запись \texttt{A} или \texttt{AAAA}.
    \item Резолвер кэширует ответ на время TTL.
    \item Резолвер возвращает ответ клиенту.
\end{enumerate}

Схема:

\[
\begin{array}{ccccc}
\text{Клиент}
& \rightarrow &
\text{Резолвер}
& \rightarrow &
\text{Корень} \\
& & \downarrow & & \\
& & \text{TLD .com} & & \\
& & \downarrow & & \\
& & \text{NS example.com} & & \\
\end{array}
\]

\subsection{Минимальный пример DNS-зоны}

\begin{lstlisting}
$ORIGIN example.com.
$TTL 3600

@       IN SOA  ns1.example.com. admin.example.com. (
            2026050601 ; serial
            3600       ; refresh
            900        ; retry
            1209600    ; expire
            3600       ; minimum
        )

@       IN NS   ns1.example.com.
@       IN NS   ns2.example.com.

ns1     IN A    192.0.2.10
ns2     IN A    192.0.2.11

www     IN A    192.0.2.20
www     IN AAAA 2001:db8::20

mail    IN A    192.0.2.30
@       IN MX   10 mail.example.com.
\end{lstlisting}

\subsection{Кэширование DNS}

Каждая DNS-запись имеет TTL --- время жизни в кэше. Если TTL равен 3600, резолвер может использовать ответ в течение одного часа, не обращаясь повторно к авторитативному серверу.

Кэширование уменьшает нагрузку и ускоряет работу, но приводит к задержке распространения изменений.

\subsection{Обратные DNS-записи}

Для обратного преобразования IP-адреса в доменное имя используются PTR-записи.

Для IPv4 применяется домен:

\[
\texttt{in-addr.arpa}.
\]

Например, для адреса:

\[
192.0.2.20
\]

обратное имя:

\[
\texttt{20.2.0.192.in-addr.arpa}.
\]

Для IPv6 используется:

\[
\texttt{ip6.arpa}.
\]

\section{Прикладные протоколы Internet}

\subsection{HTTP и HTTPS}

\begin{definition}
\textbf{HTTP} --- прикладной протокол передачи гипертекстовых документов и других ресурсов.
\end{definition}

HTTP работает по модели ``запрос--ответ''.

Пример HTTP-запроса:

\begin{lstlisting}
GET / HTTP/1.1
Host: example.com
User-Agent: ExampleClient/1.0
\end{lstlisting}

Пример HTTP-ответа:

\begin{lstlisting}
HTTP/1.1 200 OK
Content-Type: text/html
Content-Length: 42

<html><body>Hello</body></html>
\end{lstlisting}

HTTPS --- это HTTP поверх защищённого канала TLS.

\subsection{TLS}

\begin{definition}
\textbf{TLS} --- криптографический протокол, обеспечивающий конфиденциальность, целостность и аутентификацию сетевого соединения.
\end{definition}

TLS обычно используется поверх TCP. В современных системах он применяется в HTTPS, SMTP submission, IMAPS и других протоколах.

Основные задачи TLS:

\begin{itemize}
    \item согласование криптографических параметров;
    \item аутентификация сервера с помощью сертификата;
    \item выработка общих ключей;
    \item защита передаваемых данных.
\end{itemize}

\subsection{SMTP, IMAP, POP3}

\begin{itemize}
    \item \textbf{SMTP} --- отправка электронной почты.
    \item \textbf{IMAP} --- доступ к почтовому ящику с хранением писем на сервере.
    \item \textbf{POP3} --- получение почты с сервера, исторически часто с загрузкой на клиент.
\end{itemize}

\subsection{DHCP}

\begin{definition}
\textbf{DHCP} --- протокол автоматической настройки сетевых параметров узла.
\end{definition}

DHCP может выдать:

\begin{itemize}
    \item IP-адрес;
    \item маску сети;
    \item шлюз по умолчанию;
    \item DNS-серверы;
    \item время аренды адреса.
\end{itemize}

Типичный обмен DHCP:

\[
\text{Discover} \rightarrow \text{Offer} \rightarrow \text{Request} \rightarrow \text{Ack}.
\]

\section{Фрагментация и MTU}

\begin{definition}
\textbf{MTU} --- максимальный размер блока данных, который может быть передан на канальном уровне без фрагментации.
\end{definition}

Например, для Ethernet часто используется MTU 1500 байт.

\subsection{Фрагментация IPv4}

В IPv4 маршрутизатор может фрагментировать пакет, если он больше MTU следующего канала и не установлен флаг \texttt{DF}.

Фрагменты собираются только на конечном узле.

\subsection{IPv6 и Path MTU Discovery}

В IPv6 маршрутизаторы не фрагментируют пакеты. Если пакет слишком велик, маршрутизатор отбрасывает его и отправляет ICMPv6-сообщение \texttt{Packet Too Big}. Отправитель должен уменьшить размер пакетов.

\section{Безопасность в Internet}

\subsection{Основные угрозы}

\begin{itemize}
    \item подмена IP-адресов;
    \item перехват трафика;
    \item DNS spoofing;
    \item атаки отказа в обслуживании;
    \item BGP hijacking;
    \item сканирование портов;
    \item эксплуатация уязвимостей приложений.
\end{itemize}

\subsection{Механизмы защиты}

\begin{itemize}
    \item TLS для защиты прикладного трафика;
    \item DNSSEC для проверки подлинности DNS-ответов;
    \item IPsec для защиты IP-пакетов;
    \item фильтрация пакетов и межсетевые экраны;
    \item RPKI для защиты маршрутизации BGP;
    \item NAT как побочный механизм ограничения входящих соединений.
\end{itemize}

\section{Сводная схема работы Internet}

Пусть пользователь открывает:

\[
\texttt{https://www.example.com}.
\]

Полная последовательность может быть такой:

\begin{enumerate}
    \item Узел получает сетевые параметры через DHCP или статическую настройку.
    \item Браузер передаёт имя \texttt{www.example.com} системному резолверу.
    \item DNS возвращает IP-адрес.
    \item Узел определяет, находится ли адрес в локальной сети.
    \item Если адрес удалённый, пакет отправляется шлюзу по умолчанию.
    \item Для отправки кадра шлюзу определяется MAC-адрес через ARP или NDP.
    \item TCP устанавливает соединение с сервером.
    \item TLS создаёт защищённый канал.
    \item HTTP отправляет запрос.
    \item IP-маршрутизаторы пересылают пакеты по Internet.
    \item Сервер формирует HTTP-ответ.
    \item Ответные пакеты возвращаются клиенту.
    \item TCP упорядочивает байтовый поток и передаёт данные браузеру.
\end{enumerate}

\section{Краткие выводы}

\begin{enumerate}
    \item Internet является глобальной децентрализованной сетью сетей.
    \item Основой Internet является стек TCP/IP.
    \item IP отвечает за адресацию и маршрутизацию, но не гарантирует доставку.
    \item TCP обеспечивает надёжный байтовый поток, UDP --- простую передачу датаграмм.
    \item IPv4 использует 32-битные адреса, IPv6 --- 128-битные.
    \item DNS обеспечивает иерархическое сопоставление доменных имён с сетевыми данными.
    \item Маршрутизация в Internet основана на префиксах, автономных системах и протоколе BGP.
    \item Модель OSI полезна как теоретическая схема, а TCP/IP является практической архитектурой Internet.
\end{enumerate}

\end{document}
